% ============================================================================
% Report template - Sapienza Università di Roma
% Compile with: pdflatex report_template.tex (run twice for the TOC)
% Put the logo next to this file as  sapienza_logo.png  (or .pdf).
% If the logo is missing, a placeholder box is printed instead.
% ============================================================================
\documentclass[11pt,a4paper]{article}

\usepackage[utf8]{inputenc}
\usepackage[T1]{fontenc}
\usepackage[english]{babel}      % use [italian] for an Italian report
\usepackage{lmodern}
\usepackage[margin=2.5cm]{geometry}
\usepackage{graphicx}
\usepackage{xcolor}
\usepackage{amsmath,amssymb}
\usepackage{booktabs}
\usepackage{listings}
\usepackage{fancyhdr}
\usepackage{setspace}
\usepackage{caption}
\usepackage{tikz}
\usetikzlibrary{arrows.meta,positioning}
\usepackage[hidelinks]{hyperref}

% --- Colours ---------------------------------------------------------------
\definecolor{sapienza}{RGB}{124,31,52}     % adjust to the official maroon if needed
\definecolor{codebg}{RGB}{247,247,247}

% --- Metadata: edit only this block ----------------------------------------
\newcommand{\universityname}{Sapienza Università di Roma}
\newcommand{\facultyname}{Master's in Cybersecurity}
\newcommand{\coursename}{Statistics}
\newcommand{\reporttitle}{Shrimpgnature:\\ Signed Audit Logs under a FIFO Queue}
\newcommand{\reportsubtitle}{Design, implementation and performance analysis}
\newcommand{\studentnameG}{Guglielmo Forte}
\newcommand{\studentidG}{Matricola 2351773}
\newcommand{\studentnameL}{Staiano Luca}
\newcommand{\studentidL}{Matricola 2325802}
\newcommand{\professorname}{Prof. Tommaso Gastaldi}
\newcommand{\academicyear}{2026/2027}

% --- Helpers ---------------------------------------------------------------
\newcommand{\todo}[1]{\textcolor{red}{[TODO: #1]}}   % delete all before submitting

\lstdefinelanguage{JavaScript}{
  keywords={const,let,var,function,async,await,return,if,else,for,while,new,
            export,import,try,catch,throw,true,false,null,typeof},
  comment=[l]{//}, morecomment=[s]{/*}{*/}, morestring=[b]', morestring=[b]",
  morestring=[b]`
}
\lstset{
  language=JavaScript, basicstyle=\ttfamily\footnotesize,
  backgroundcolor=\color{codebg}, keywordstyle=\color{sapienza}\bfseries,
  commentstyle=\color{gray}\itshape, stringstyle=\color{blue!60!black},
  frame=single, rulecolor=\color{black!20}, breaklines=true,
  numbers=left, numberstyle=\tiny\color{gray}, xleftmargin=1.5em,
  framexleftmargin=1.5em, columns=fullflexible
}

% --- Header / footer -------------------------------------------------------
\pagestyle{fancy}
\fancyhf{}
\lhead{\small\textcolor{sapienza}{\universityname}}
\rhead{\small\coursename}
\cfoot{\thepage}
\renewcommand{\headrulewidth}{0.4pt}
\setlength{\headheight}{14pt}

\begin{document}

% ============================== TITLE PAGE ==================================
\begin{titlepage}
\centering

% --- LOGO LEFT ALIGNMENT START ---
\begin{flushleft}
\IfFileExists{logo_Sapienza.png}{%
    \includegraphics[width=0.33\textwidth]{./logo_Sapienza.png}%
}{\IfFileExists{sapienza_logo.pdf}{%
    \includegraphics[width=0.33\textwidth]{}%
}{%
    \fbox{\parbox[c][3.5cm][c]{5cm}{\centering\textcolor{red}{LOGO MISSING}\\ \small add \texttt{sapienza\_logo.png}}}%
}}
\end{flushleft}
\nopagebreak[4]
\vspace{1.2cm}
% --- LOGO LEFT ALIGNMENT END ---

{\Large\scshape \universityname\par}
\vspace{0.2cm}
{\large \facultyname\par}
{\large \coursename\par}
\vspace{1 cm}
{\color{sapienza}\rule{\textwidth}{1pt}}\\[0.6cm]
{\Huge\bfseries \reporttitle\par}
\vspace{0.4cm}
{\Large \reportsubtitle\par}
\vspace{0.3cm}
{\color{sapienza}\rule{\textwidth}{1pt}}
\vfill
\begin{tabular}{rl}
\textbf{Students:} & \studentnameG \\
& \studentidG \\[0.3cm]
& \studentnameL \\
& \studentidL \\[1cm]
\textbf{Professor:} & \professorname \\[3cm]
\end{tabular}
\vfill
{\large Academic Year \academicyear\par}
\end{titlepage}


% ============================== FRONT MATTER ================================
\pagenumbering{roman}
\begin{abstract}
\noindent
This report presents the ? 
\end{abstract}
\tableofcontents
\clearpage
\pagenumbering{arabic}

\begin{spacing}{1.1}

% ============================== 1. INTRODUCTION (~1 page) ===================
\section{Introduction}

Digital signatures are essential for proving the authenticity, integrity and non-repudiation of documents. They guarantee that a specific record was processed by a notary and that it has not sustained any modification since being signed. \\\\
However if a system solely relied on digital signatures, a malicious actor could silently delete a previously notarized document, reorder historycal transactions or present an incomplete version of the events to an auditor which would not detect any related issues since the signatures of the individual document would still be mathematically valid.\\\\
To prevent this, a notary requires a notary requires an ordered log which enables tamper detection. By structuring recepits into a caryptographic hash chain, each new notary operation is linked to the previous one, thus creating an immutable sequence. Any attempt to remove, reorder, alter past events will inevitably cause incoherence in the chain, which is detectable by the system. Thus \textbf{digital signatures} are needed to ensure fine-grained validity of a document, while the \textbf{hash chain} operates at a coarse-grained level.

\subsection{Goals}
The main goal of this project is to simulate and analize a central notary authority that processes document signing requests. The system includes a performance simulation component that models the signing authority as a single-server queue, allowing for empirical observation of queueing metrics and comparison against theoretical models such as M/M/1 and M/G/1.

\subsection{Structure of the report}
The report is structured as follows: 
\begin{itemize}
\item{\textbf{Section 2} provides background informations on digital signatures, hash chains and queueing models.}
\item{\textbf{Section 3} outlines the system architecture, describing components, message flows and data structures.}
\item {\textbf{Section 4} details the software implementation of the queue, the worker processes, the verification entity and the traffic generator.}
\item {\textbf{Section 5} presents an experimental evaluation of service and waiting times under load.}
\item {\textbf{Section 6} epitomizes the work and concludes the report.}
\end{itemize}
% ============================== 2. BACKGROUND (~2 pages) ====================
\pagebreak

\section{Background}

\subsection{Digital signatures}
Digital signatures ensure the authenticity, non-repudiation and integrity of digital documents. The central notary authority supports two asymmetric cryptographic algorithms: the Elliptic Curve Digital Signature Algorithm (\textit{ECDSA}) using the P-256 curve, and the RSA algorithm with the Probabilistic Signature Scheme (\textit{RSA-PSS}) using a 4096-bit modulus.
\\When a document digest is submitted, the authority signs it with its private key, generating a signature that is verifiable by anyone using the corresponding public key.

\subsection{Hash chains and tamper evidence}
To ensure the immutability of the \emph{audit log}, the system structures recepits into a \textit{hash chain}. Each new recepit computes a \texttt{chainHash} using the SHA-256 hash function over the previous block's hash \texttt{prevHash}, the current block's \texttt{signature} and the \texttt{docDigest}; thus forming a sequential dependency which guarantees the propagation of any tampering attempt made on any past block invalidating every subsequent one.
\\\\
The logical structure of the afore mentioned blocks can be visualized as follows:
\begin{equation}
  h_i = \mathrm{SHA\text{-}256}\bigl(h_{i-1} \,\Vert\, \sigma_i \,\Vert\, d_i\bigr),
  \qquad h_0 = 0^{256},
  \label{eq:chain}
\end{equation}
where $d_i$ is the document digest and $\sigma_i$ its signature.
\\\\
If an attacker modifies the contents of entry $j$, for example altering the document digest, the output of the hash function for that block changes entirely due to the avalanche effect, resulting in a new hash $h'_j$ \neq $h_j$. Because the next block's hash, $h_{(j+1)}$ depends on $h_j$, computing $H(h'_j \parallel \text{data}_{j+1})$ will yield a different result, $h'_{(j+1)}$. By induction this cascades forward indefinitely.

\pagebreak

\subsection{Queueing model}
The system evaluates performance using queueing theory, modelling the authority as a single-server queue with Poisson arrivals. Two analytical models are tracked to predict queue waiting times ($W_q$):

\begin{itemize}
  \item \textbf{M/M/1}: assumes both exponential inter-arrival times and exponential service times. Its infinitesimal generator matrix $Q$ is:
  \[
  Q = \begin{pmatrix}
  -\lambda & \lambda & 0 & 0 & \cdots \\
  \mu & -(\lambda + \mu) & \lambda & 0 & \cdots \\
  0 & \mu & -(\lambda + \mu) & \lambda & \cdots \\
  0 & 0 & \mu & -(\lambda + \mu) & \cdots \\
  \vdots & \vdots & \vdots & \vdots & \ddots
  \end{pmatrix}
  \]
  The expected waiting time in the queue is $W_q = \frac{\rho}{\mu (1 - \rho)}$.
  
  \item \textbf{M/G/1}: assumes exponential arrivals but general service times. It uses the Pollaczek-Khinchine formula: $W_q = \frac{\lambda(\sigma_S^2 + \bar{S}^2)}{2 (1-\rho)}$.
  Here, $\lambda$ is the arrival rate, $\mu$ is the service rate, $\rho = \frac{\lambda}{\mu}$ is the system load, $\bar{S}$ is the mean service time, and $\sigma_S^2$ is the service time variance.
\end{itemize}

\pagebreak

\section{System Architecture}

\subsection{Components}

The system consists of multiple modules operating asynchronously across the main browser thread and a Web Worker:

\begin{itemize}
\item \textbf{UI and Telemetry View}: handles DOM manipulations, dynamic charting with \texttt{Chart.js}, and real-time KPI updates.
\item \textbf{Traffic Generator}: produces synthetic document signing requests according to a Possion process.
\item \textbf{Central Notary Authority}: a Web Worker responisble for cryptographic key generation, document signing and hash chain computation.
\item \textbf{Verification Engine}: audits the public ledger to detect tampering or history-rewriting attacks.
\end{itemize}


\subsection{Message flow}

When a request is generated, its payload (a document digest and algorithm choice) is enqueued in the FIFO Server. Once the server is idle, it dispatches the request via a \texttt{SIGN\_SUCCESS} message containing a receipt and telemetry data. The main thread then handles this message by updating the public audit log, telemetry metrics and the UI.

\subsection{Receipts and telemetry}
The generated receipt includes the sequence number, document digest, generated signature, chosen algorithm, timestamp and a \texttt{chainProof} (containing \texttt{prevHash} and \texttt{chainHash}). Parallel to this, telemetry data is captured for performance monitoring, tracking specifically \texttt{tArrival} (time of queue entry), \texttt{tQueueExit} (time of dispatch to the worker) and \texttt{serviceTimeSeconds} (duration of cryptographic operations).

\section{Implementation}
\todo{mettere codice per ste implementazioni}

\subsection{FIFO queue and single-server state machine}
The queue logic is implemented in \texttt{fifoServer.js}. It encapsulates an array-based queue and maintains the server state as either \texttt{IDLE} or \texttt{PROCESSING}. When a job is enqueued, it receives an arrival timestamp. If the server in \texttt{IDLE}, the job is immediately dequeued, its queue exit time is stamped, the state transitions to \texttt{PROCESSING}, and the request is dispatched. Upon task completion, the state resets to \texttt{IDLE}, and the \texttt{next()} routine is invoked.

\subsection{Signing and hash chaining in the worker}
The worker initializes by generating ECDSA and RSA-PSS key pairs using the native \texttt{crypto.subtle} API. During processing, the signing method is invoked over the document digest. Immediately following the signature generation, the hash chain link is computed by calling \texttt{computeHashChain}. Performance metrics strictly measure the time around these cryptographic bounds. Executing this inside a Web Worker prevents CPU-intensive 4096-bit RSA operations from blocking the main thread.


\subsection{Verification engine}
The audit engine (\texttt{auditLedger.js}) performs a two-layer verification sequentially over the ledger. The first one involves recalculating the \texttt{chainHash} for each block and ensuring it matches the \texttt{prevHash} of the subsequent block. The second layer imports the Notary Authority's public key and verifies every signature. The engine reports the sequence number (\texttt{firsTamperedBlock}or \texttt{firstForgedSignatureBlock}) where any integrity check fails.


\subsection{Traffic generator}
The synthetic workload is managed by \texttt{poissonGenerator.js}. It uses exponential inter-arrival times derived via the inverse transform method: delay $= (-\frac{\ln(1-\text{rand}())} {\lambda}) \times 1000$ milliseconds. A recursive timeout function schedules subsequent submissions, allowing to continously simulate requests matching the user-configured arrival late ($\lambda$)








\section{Experimental Evaluation}

\subsection{Setup}
\todo{Browser and version, CPU, OS, number of requests per run, number of repetitions.
Mention that service time is measured inside the worker and covers signing plus the
chain hash, not message passing.}

\subsection{Service time of ECDSA vs RSA-PSS}
\begin{table}[ht]
\centering
\begin{tabular}{lrrrr}
\toprule
Algorithm & $\mathbb{E}[S]$ (ms) & $\mathrm{Var}(S)$ & $\mathbb{E}[S^2]$ & $\mu=1/\mathbb{E}[S]$ (req/s)\\
\midrule
ECDSA P-256   & \todo{} & \todo{} & \todo{} & \todo{} \\
RSA-PSS 4096  & \todo{} & \todo{} & \todo{} & \todo{} \\
\bottomrule
\end{tabular}
\caption{Measured service-time statistics.}
\end{table}

\subsection{Waiting time versus load}
\todo{For $\rho \in \{0.3, 0.5, 0.7, 0.9\}$ choose $\lambda=\rho/\mathbb{E}[S]$, run the
generator, and compare the measured mean wait with Eq.~\eqref{eq:pk}. Plot
measured vs.\ predicted (one curve per algorithm).}

\begin{table}[ht]
\centering
\begin{tabular}{lrrrr}
\toprule
Algorithm & $\rho$ & $\lambda$ (req/s) & $\overline{W_q}$ measured (ms) & $\mathbb{E}[W_q]$ P--K (ms)\\
\midrule
\todo{} & \todo{} & \todo{} & \todo{} & \todo{} \\
\bottomrule
\end{tabular}
\caption{Measured vs.\ predicted mean waiting time.}
\end{table}

\subsection{Discussion}
\todo{Does the model fit? Possible sources of error: \texttt{setTimeout} jitter, short runs,
warm-up (first requests are slower), service times not independent, other browser load.}


% ============================== 7. CONCLUSION (~0.5 page) ===================
\section{Conclusion}
\todo{Summarise results in 3--4 sentences, state one limitation, and give two
concrete future-work items (e.g.\ multiple workers $\to$ M/G/c, persistent log,
external timestamping).}

\end{spacing}

% ============================== REFERENCES ==================================
\begin{thebibliography}{9}
\bibitem{haber}  S.~Haber, W.~S.~Stornetta. How to time-stamp a digital document.
  \emph{Journal of Cryptology}, 3(2), 1991.
\bibitem{fips}   NIST. \emph{Digital Signature Standard (DSS)}, FIPS 186-5, 2023.
\bibitem{rfc}    K.~Moriarty et al. \emph{PKCS \#1: RSA Cryptography Specifications
  Version 2.2}, RFC 8017, 2016.
\bibitem{kleinrock} L.~Kleinrock. \emph{Queueing Systems, Volume 1: Theory}. Wiley, 1975.
\bibitem{webcrypto} W3C. \emph{Web Cryptography API}. \url{https://www.w3.org/TR/WebCryptoAPI/}
\end{thebibliography}

\end{document}